\input{../tex-inputs/latex-leseansicht-vorspann}

\section[Gustav Schwarzkopf an Arthur Schnitzler, 17. 8. 1899]{L04295 Gustav Schwarzkopf an Arthur Schnitzler, 17. 8. 1899}
\nopagebreak\mylabel{L04295v}
\rehead{ }\normalsize\beginnumbering\briefempfaengerindex{Schnitzler, Arthur@\textsc{Schnitzler, Arthur}!zzzSchwarzkopf, Gustav@\emph{von Gustav Schwarzkopf}!1899-08-173@{17. 8. 1899}|(be}
\toendnotes[C]{\smallbreak\pagebreak[2]}
\correspDesc{Versand  durch Gustav Schwarzkopf am 17. 8. 1899 in Wien
\newline{}Erhalt  durch Arthur Schnitzler im Zeitraum [18. 8. 1899 – 22. 8. 1899?] in Bad Ischl}\toendnotes[C]{\smallbreak}
\Standort{CUL, Schnitzler, B 96a.}
\physDesc{Brief, 1 Blatt, 4 Seiten, 2281 Zeichen
\newline{}Handschrift: schwarze Tinte, deutsche Kurrent}\toendnotes[C]{\smallbreak}
\pstart
           \raggedleft{}{\pb}17. 8. 99\pend
           \vspace{0.5em}
\pstart
           Lieber Arthur, die beiden Anſichtskarten \label{K_L04295-1v}\edtext{aus \textsc{Predazzo\oindex{Predazzo@\textbf{Predazzo}, \emph{Verwaltungsgebiet}|pw}}}{\lemma{\textnormal{\emph{aus Predazzo}}}\Cendnote{\textnormal{{XXXX ref} XXXX}}}\label{K_L04295-1} und \textsc{\label{K_L04295-2v}\edtext{Martino\oindex{San Martino di Castrozza@\textbf{San Martino di Castrozza}|pw}}{\lemma{\textnormal{\emph{Martino}}}\Cendnote{\textnormal{{XXXX ref} XXXX}}}\label{K_L04295-2}} habe ich \textsc{Samstag} reſpek. So{\geminationn}tag, alſo erſt nach
               Abſendung meiner \label{K_L04295-3v}\edtext{Karte}{\lemma{\textnormal{\emph{Karte}}}\Cendnote{\textnormal{{XXXX ref}
                  XXXX}}}\label{K_L04295-3} erhalten. Warum{ }ſind Sie nicht länger im \textsc{Martino\oindex{San Martino di Castrozza@\textbf{San Martino di Castrozza}|pw}} geblieben? Hat es Ihnen nicht zugeſagt? Aus der Tatſache, daß Sie einen \uline{Brief} von mir beko{\geminationm}en –
               um meine Ankunft anzuzeigen, hätte ja eine Karte genügt – werden Sie{ }ſchon erraten
               haben, daß ich nicht die Abſicht habe nach \textsc{Ischl\oindex{Bad Ischl@\textbf{Bad Ischl}|pw}} zu ko{\geminationm}en. Meine Trägheit, die Unluſt etwas zu
                  begi{\geminationn}en oder gar eine Veränderung mit mir{ }ſelbſt
               vorzunehmen, iſt{ }ſo groß geworden, daß ich die notwendige Energie {\pb}nicht aufbringen ka{\geminationn}. Das klingt unwahrſcheinlich, wie eine Ausrede, aber
               es iſt wirklich{ }ſo, glaub mir. Selbſtverſtändlich ka{\geminationn}
               man{ }ſich von dieſer Apathie durch einen Ruck befreien, und das würde ich wol auch
               tun, we{\geminationn} ich nicht wüßte, daß Sie in vier Wochen ja doch
               wieder in Wien\oindex{Wien@\textbf{Wien}, \emph{Verwaltungsgebiet}|pw}{ }ſein werden. Jetzt würden Sie
               übrigens einen{ }ſchlechten Geſellſchafter an mir haben, – Sie wiſſen ja{ }ſchon, was das
               bei mir heißt – der Ihre »mäßige Sti{\geminationm}ung« nur
               verſchlechtern kö{\geminationn}te. Ich bin gründlich verſti{\geminationm}t,– ganz wörtlich zu nehmen – {\pb}wie ein Inſtrument, das lange nicht
               geſpielt wurde. Seit einem Monat hab ich ja{ }ſo{ }ſelten geſprochen – nur hie und da mit
                  \textsc{Eberma{\geminationn}\pwindex{Ebermann, Leo 16.\,7.\,1863 Draganovka – 9.\,10.\,1914 Wien@\textsc{Ebermann, Leo} (16.\,7.\,1863 Draganovka – 9.\,10.\,1914 Wien), \emph{Schriftsteller, Journalist, Rechtswissenschaftler}|pw}} und \textsc{Kapper\pwindex{Kapper, Hans 10.\,3.\,1891 Wien – September 1962 New York City@\textsc{Kapper, Hans} (10.\,3.\,1891 Wien – September 1962 New York City), \emph{Schriftsteller}|pw}} – und dafür{ }ſoviel Unnützes und Unerfreuliches gedacht. – \textsc{Eberma{\geminationn}\pwindex{Ebermann, Leo 16.\,7.\,1863 Draganovka – 9.\,10.\,1914 Wien@\textsc{Ebermann, Leo} (16.\,7.\,1863 Draganovka – 9.\,10.\,1914 Wien), \emph{Schriftsteller, Journalist, Rechtswissenschaftler}|pw}} reiſt Ende dieſer Woche fort nach Salzburg\oindex{Salzburg@\textbf{Salzburg}, \emph{Verwaltungsgebiet}|pw},
               wird da{\geminationn} auch nach \textsc{Ischl\oindex{Bad Ischl@\textbf{Bad Ischl}|pw}} ko{\geminationm}en. – »Denk mal an«, vorgeſtern habe ich hier
                  \textsc{Adele Sandrock\pwindex{Sandrock, Adele 19.\,8.\,1863 Rotterdam – 30.\,8.\,1937 Berlin@\textsc{Sandrock, Adele} (19.\,8.\,1863 Rotterdam – 30.\,8.\,1937 Berlin), \emph{Schauspielerin}|pw}} geſprochen. In meiner Straße, auf dem »tiefen
                  Graben\oindex{Wien@\textbf{Wien}!I., Innere Stadt@\textbf{I., Innere Stadt}!Tiefer Graben@\textbf{Tiefer Graben}, \emph{Straße}|pw}«. Sie hatte die Gnade mich anzuſprechen, ich hätte es gewiß nicht
               getan;{ }ſie{ }ſieht gut aus, wird{ }ſtark, was{ }ſie{ }ſehr beſorgt macht, {\pb}und erzählte mir, daß{ }ſie 180 mal
               geſpielt und alle Stücke{ }ſelbſt inſcenirt \substVorne{}\textsuperscript{\textcolor{gray}{hen}}\substDazwischen{}und\substHinten{} glänzende Geſchäfte gemacht habe. »In meinem Koffer iſt jetzt der Ehrgeiz
               tief unten verpackt und das Checkbuch ganz oben«{ }ſagte{ }ſie. Buchbinder\pwindex{Buchbinder, Bernhard 7.\,7.\,1849 Budapest – 24.\,6.\,1922 Wien@\textsc{Buchbinder, Bernhard} (7.\,7.\,1849 Budapest – 24.\,6.\,1922 Wien), \emph{Schriftsteller, Journalist, Schauspieler}|pw} kö{\geminationn}te aus dieſer
               Begegnung zwei reizende Feuilletons machen. Als wir uns verabſchiedet hatten, dachte
               ich mir: Gut iſt es nur, daß sie gewiß längſt meinen Namen vergeſſen hat,{ }ſonſt würde{ }ſie doch{ }ſicher erzählen: »Neulich habe ich den \textsc{Schw.},
               dieſen aus dem Zuchthaus entſprungenen Idioten getroffen –.« – – –\pend
           
\pstart
           Bitte, empfehlen Sie mich allen Mitgliedern Ihrer Familie und grüßen Sie auch Ihre
               lieben kleine Neffen\pwindex{Schnitzler, Hans 11.\,7.\,1895 Wien – 26.\,3.\,1967 Chicago@\textsc{Schnitzler, Hans} (11.\,7.\,1895 Wien – 26.\,3.\,1967 Chicago), \emph{Chirurg}|pwv}
               herzlichſt vom »Onkel Glatzkopf«.\pend
           
\pstart
           Herzlich{\\[\baselineskip]}Ihr{\\[\baselineskip]}\spacefill\mbox{GustavSch}\pend
           \leftskip=0em{}\selectlanguage{ngerman}\endnumbering\briefempfaengerindex{Schnitzler, Arthur@\textsc{Schnitzler, Arthur}!zzzSchwarzkopf, Gustav@\emph{von Gustav Schwarzkopf}!1899-08-173@{17. 8. 1899}|)be}\mylabel{L04295h}
\begin{anhang}
\end{anhang}\newcommand{\dateiname}{L04295}\newcommand{\titel}{Gustav Schwarzkopf an Arthur Schnitzler, 17. 8. 1899}\newcommand{\editorInnen}{Herausgegeben von Jahnke, SelmaMüller, Martin Anton}\input{../tex-inputs/latex-leseansicht-abspann}
